\documentclass[12pt, a4paper, oneside]{report}
\usepackage[spanish]{babel}
\usepackage[utf8x]{inputenc}
\usepackage[T1]{fontenc}
\usepackage{newtxtext}
\usepackage{etoolbox}
\usepackage[autostyle, spanish=spanish]{csquotes}
\usepackage{amsmath,amssymb,amsfonts}
\usepackage[margin=2.5cm]{geometry}
\usepackage{setspace}
\usepackage{graphicx}

\usepackage[table]{xcolor}

\definecolor{negro}{HTML}{000000} % #000000
\definecolor{rojo}{HTML}{ff0000} % #ff0000
\definecolor{verde}{HTML}{2da44e} % #2da44e

\graphicspath{{./img/}}

\title{Dosier de prácticas de los Temas 3 y 4}
\author{González Molina, José Ángel}
\date{\today}
\setlength{\parskip}{1em}
\setlength{\parindent}{1cm}
\onehalfspacing

\begin{document}
    \color{negro}
    \maketitle
    \clearpage

    \renewcommand{\chaptername}{Ejercicio}
    \renewcommand{\figurename}{Fig.}
    \chapter{Elabora un texto descriptivo}
        \begin{figure}[ht]
            \centering
            \includegraphics[scale = .75]{balthus}
            \caption{Balthus (1938). Thérèse soñando.}
        \end{figure}
        Al entrar en la habitación se encontró con Thérèse dormitando. La iluminación era tenue debido a que
        todas las lámparas se encontraban apagadas por lo que poco a poco se iban vislumbrando más detalles de
        la estancia.

        Una panorámica de la escena mostraba una toalla que se encontraba descuidadamente depositada en la
        mesa. Frente a esta, una silla mal colocada. En el centro se podía observar a la pequeña recostada en
        un cojín esmeralda, el cual recogía su espalda. Finalmente en el suelo junto a ella un felino níveo
        como la leche, ajeno en apariencia al sopor de su compañera de juegos. Aquel día Thérèse estaba
        vestida con una camisa blanca y su habitual falda escarlata. Apoyada totalmente en el respaldo y con
        los brazos detrás de su cabeza, ojos cerrados y semblante apaciguado.

        Parecía que las tareas domésticas habían sido interrumpidas por el cansancio, o al menos eso decía el
        estado desordenado de la habitación. Pensó en despertarla, pero si no lo había hecho el ruido que
        estaba haciendo el gato con el plato, era porque el descanso era merecido.
        \clearpage
    \chapter{Elabora una introducción breve de un trabajo cuyo índice es el siguiente:}
        \section*{Índice}
        \begin{Large}
            \renewcommand{\labelenumii}{\arabic{enumi}.\arabic{enumii}}
            \begin{enumerate}
                \item Los orígenes del hombre
                \item Antes del \emph{Homo sapiens}. Los colonizadores del Viejo Mundo
                \item Los neandertales: ¿los primeros hombres modernos?
                \begin{enumerate}
                    \item Los primeros artistas
                    \item Pastores y agricultores
                \end{enumerate}
                \item Conclusiones
                \item Bibliografía
            \end{enumerate}
        \end{Large}
        \clearpage
        \section*{Introducción}
            En el presente trabajo se investigarán las poblaciones humanas en Europa en el periodo de tiempo
            anterior al apogeo del Homo sapiens. El \emph{Homo sapiens} es la especie a la que pertenecen
            todos los seres humanos que viven en la actualidad. Como se verá existieron multitud de homínidos
            \emph{hominina}, es decir, la familia de primates bípedos a la que pertenecemos y que contiene a
            otras especies que también poblaron Europa como el \emph{Homo antecessor}, el
            \emph{Homo heidelbergensis} y el \emph{Homo neanderthalensis}. Nos centraremos principalmente en
            los neandertales y en cómo colonizaron Europa llegando a compartir periodo temporal con nuestra
            especie.

            A lo largo de la historia de la antropología y la paleontología, ha existido una dificultad clara
            a la hora de catalogar con exactitud las especies y momentos temporales implicados en la
            colonización del continente europeo. Una de las tesis que más aceptación tiene es la capacidad de
            creación de arte o el empleo de técnicas para seleccionar animales a convertir en ganado o
            semillas que cultivar. Pero ¿son los \emph{Homo sapiens} los primeros en conseguir estos logros?

            Como se expondrá en el capítulo sobre los neandertales, las nuevas evidencias y hallazgos nos
            inducen a pensar en que dichos logros fueron logrados con anterioridad. En un primer apartado se
            describirán las últimas muestras de arte rupestre, datadas en fechas muy anteriores a los
            registros que se conservan de los primeros sapiens. Además se ha realizado un análisis de la
            complejidad y variedad de dichas obras en una comparativa a aquellas que ya conocíamos y de origen
            más reciente.

            Seguidamente, se abordará el estudio de las muestras halladas en recientes excavaciones donde se
            encontraron herramientas compatibles con un cultivo complejo de semillas. Dichas excavaciones
            están realizadas en terrenos con datación geológica muy anterior a los periodos habituales donde
            se sitúan a los sapiens. También se analizarán los resultados aportados por los restos fósiles
            hallados en Irak donde se encontraron trazas de alimentos cocinados.

            Finalmente el trabajo concluye con las conclusiones a las que nos pueden llevar las
            investigaciones descritas con anterioridad y las consecuencias que pueden tener en nuestro
            conocimiento del origen del ser humano y como ha ido evolucionando hasta la especie dominante del
            mundo de la que formamos parte ahora.
        \clearpage

        \chapter{Elabora un breve estado de la cuestión a partir de la selección de textos que se presentan a continuación}
            \setcounter{section}{0}
            \section{Textos que se deben utilizar}
                \subsection{Pérez, Antonio
                El lenguaje de los homínidos, 2009
                Madrid, Cátedra.}

                El lenguaje hablado es, hoy por hoy, el único rasgo cualitativo que diferencia drásticamente
                a la humanidad del resto de los animales. La complejidad y funcionalidad de la comunicación
                oral codificada entre los humanos parece estar fuera del alcance de cualquier otra criatura
                del planeta, con excepción quizá de los cantos de las ballenas corcovadas, que parecen tener
                connotaciones gramaticales. (p. 11)

                \subsection{Correas, María
                “Origen del lenguaje humano”, 2007
                Revista Cuestiones científicas, 35, pp. 1-15}

        El lenguaje surge de forma abrupta y su emergencia está ligada a la aparición de Homo sapiens, hace 40.000 años en Eurasia. No existió un protolenguaje del que las lenguas modernas evolucionaran. La aparente explosión simbólica del Paleolítico Superior europeo iría acompañada de un aumento drástico en la complejidad y estructuración lingüística (p. 12).
        \subsection{Campos, Andrea
        “Homínidos”, 2003
        Capítulo del libro ed. por Asunción Rayo, Estudios de lingüística y prehistoria
        Cádiz, Universidad de Cádiz, pp. 123-155}
        Los neandertales y demás homínidos pre sapiens poseían un lenguaje rudimentario, muy lejos del lenguaje tal y como lo conocemos. Las ventajas adaptativas del lenguaje moderno serían suficientes para explicar los cambios conductuales que supuestamente acaecieron durante la transición Paleolítico Medio-Superior (pp. 124-125).
        \subsection{Gutiérrez, Salvador
        “Encefalización y evolución cultural”, 2004
        Revista Lengua y Evolución, número 7, pp. 24-67}
        El lenguaje evoluciona gradualmente a partir de las primeras poblaciones de Homo ergaster o erectus que comenzaron a salir de África. El estudio de cráneos y réplicas internas del cerebro de Homo habilis o de Homo erectus de hace 1,55 millones de años ha permitido afirmar que estas estructuras apenas han cambiado desde entonces y que, por lo tanto, los primeros Homo gozaban de la capacidad neurológica necesaria para producir y emplear un lenguaje articulado (p. 50).
        \subsection{Sisti, Ana María
        Lenguaje humano: cuándo y por qué, 2009
        Roma, Lípari}
        En la línea continuista existe un modelo derivado de la morfología y etología comparada de primates que sugiere que la necesidad de grupos mayores entre nuestros primeros ancestros fue la que condujo a la evolución del lenguaje y a la encefalización de los homínidos. La estrecha relación entre el tamaño de los grupos y la encefalización permite situar el punto en la evolución de los homínidos en el cual el tamaño del grupo se hizo lo suficientemente grande como para que el lenguaje se hiciese necesario a la hora de mantener la cohesión social. Una conclusión preliminar es que el lenguaje surgió como una forma de lazo de conexión más efectivo en el empleo del tiempo social. El tamaño del grupo está limitado por el número de relaciones que un individuo puede controlar exitosamente, número que a su vez está limitado por el tamaño relativo del neocórtex (pp. 114-15).
        \section*{Estado de la cuestión}
        \clearpage

    \chapter{Prácticas de citación}
        \section{Señala y corrige los errores de uso del sistema APA}
        \begin{Large}
            \renewcommand{\labelenumi}{\alph{enumi})}
            \begin{enumerate}
                \item Según señala Gómez (1988: pp. 24-26), la biodiversidad de las zonas tropicales está en
                claro retroceso.

                \textit{Es una cita indirecta narrativa, por lo que el año debe aparecer junto al autor y las
                páginas al final de la cita, siguiendo el formato APA adecuado:}

                Según señala Gómez (1988), la biodiversidad de las zonas tropicales está en claro retroceso
                (pp. 24-26).
                \clearpage

                \item Si comparamos los resultados de las distintas investigaciones, llegaremos a la
                conclusión de que las zonas tropicales, frente a otros ecosistemas, están en claro retroceso
                (Gómez 1988).

                \textit{Es una cita indirecta parentética, por lo que la cita debe aparecer al final de la
                frase parafraseada, en el formato adecuado}

                Si comparamos los resultados de las distintas investigaciones, llegaremos a la conclusión de
                que las zonas tropicales, frente a otros ecosistemas, están en claro retroceso (Gómez, 1988).
                \item Aparicio (2013) (Linguistics, n.º 12, p. 124) señala que ``las zonas tropicales, frente
                a otros ecosistemas, están en claro retroceso''.

                \textit{Esta cita es directa y narrativa por lo que la norma es mencionar el año junto al
                autor y la página con el formato adecuado al final de la cita. El resto de información
                aparecerá en la entrada bibliográfica.}

                Aparicio (2013) señala que “las zonas tropicales, frente a otros ecosistemas, están en claro
                retroceso” (p. 124).
                \clearpage
                \item Si comparamos los resultados de las distintas investigaciones, llegaremos a la
                conclusión de que las zonas tropicales, frente a otros ecosistemas, están en claro retroceso.
                (Gómez, 1988:3, Alvarado y Suñer, 1999:29, Carrasco, 2000:11)

                \textit{Esta cita es indirecta y parentética, pero de varios autores que coinciden en la idea.
                Se deben citar los autores al final de la frase en orden alfabético y con el formato
                correcto:}

                Si comparamos los resultados de las distintas investigaciones, llegaremos a la conclusión de
                que las zonas tropicales, frente a otros ecosistemas, están en claro retroceso
                (Alvarado y Suñer, 1999, p. 29; Carrasco, 2000, p. 11; Gómez, 1988, p. 3).
            \end{enumerate}
        \end{Large}

        \section{Elabora una cita directa (con tres líneas y con todas las líneas) a partir de la siguiente
        información:}

            Artículo de David Gil de 2008 titulado "¿Cuán complejas son las lenguas aislantes?", de la revista Language, número 21, pp. 109-132.

            Las lenguas pueden ser más o menos complejas en virtud del número de elementos que posean. Se consideran lenguas complejas a aquellas que tienen un número alto de palabras, emplean morfemas para determinadas relaciones sintácticas y tienen más de diez fonemas vocálicos y veinticinco consonánticos. Las lenguas que tienen menos de 2.000 palabras, carecen de morfemas y poseen un número reducido de fonemas son lenguas simples.

        \section{Elabora una cita indirecta (una parentética y otra narrativa) a partir de la siguiente información:}

            Libro de Johanna Nichols titulado Diversidad lingüística en el espacio y en el tiempo, publicado en 1992 en la editorial University of Chicago Press (en Chicago).

            Existen ciertas áreas del planeta que condensan un número considerable de lenguas complejas. Parece, pues, que la complejidad lingüística se circunscribe a ciertas regiones. Esto significa que la complejidad de las lenguas puede estar relacionada con una geografía determinada, como las zonas aisladas (p. ej. montañas, selvas tropicales, etc.).

\end{document}
